\documentclass[11pt,a4paper]{article}
\usepackage[T1]{fontenc}
\usepackage{lmodern,microtype,amsmath,amssymb,enumitem}
\usepackage[margin=1in]{geometry}
\usepackage[hidelinks]{hyperref}
\setlength{\parindent}{0pt}
\setlength{\parskip}{5pt}
\begin{document}
\begin{center}
{\Large\bfseries Staking and Rotation Module: Review of v1}\\[4pt]
10 October 2026
\end{center}
\textbf{Frozen source:} \texttt{formalization/staking\_rotation\_module\_v1.tex}.
\textbf{SHA-256:}\\
\texttt{7f033b35432940afc9f8799c2c50b81fc0e39bf6341d5ef4da4d1ebf00bdcb64}.

\section*{Summary}
\textbf{Final verdict.} Do not accept v1 as a complete specification.
The conditional prefix-safety induction is sound under the explicit
regional BFT, terminal-closure and finalized-L1 contracts. The broader
progress and churn statements contain critical errors, and some transition
branches are incomplete. These findings do not demonstrate a fork under
the stated closure contract.

This record consolidates three independent interface/security and stake/lifecycle
verifier reports. The design source remained frozen during verification.
The installation check was also identified by the author. Repairs
belong in v2, not in the reviewed source.

\section*{Findings and verification log}
\begin{enumerate}[label=\textbf{V1-\arabic*.},leftmargin=*,itemsep=9pt]
\item \textbf{Critical Error: opening progress lacks eligibility.}
Lines 357--359 state that eventual L1 actions, seed publication, data
recovery and closure let the module ``open an eligible region.'' All those
premises can hold while the status map keeps every region off, or every
region is empty. A new configuration may install idle; production cannot
start. The progress conclusion needs a separate eligibility/service premise.

\item \textbf{Justification Gap: idle resumption is inconsistent.}
Lines 232--233 select the first eligible entry ``from idle.'' R4,
lines 281--282, instead sends every non-epoch-change transition to the
within-epoch successor rule. That rule needs a previous region, but
the current session is $\bot$ and no previous-region cursor is stored.
The intended first-eligible branch must be made explicit.

\item \textbf{Critical Error: the gross-change bound includes unbounded
penalty effects.}
Lines 124--125 say the budgets bound ``total changes''; lines 159--161
rebase snapshots to remaining cutoff balances. With seven 10-unit
positions, an outbound budget of 10, no EXIT, and two 9-unit slashes,
the next total falls from 70 to 52. The loss of 18 exceeds the budget.
The two affected positions can have old Byzantine weight $20<70/3$,
so this counterexample need not violate the regional fault assumption.
The budget can bound voluntary queued changes, not arbitrary penalty losses.

\item \textbf{Justification Gap: zero balances have no membership rule.}
The table requires positive weights at lines 55--56. Applying queued
exits/joins and recording cutoff balances leaves a fully debited incumbent
with zero weight when it did not request exit. Omission, old-weight
retention, and zero-weight retention are different transitions. Automatic
removal and its retirement accounting must be specified.

\item \textbf{Justification Gap: canceled FIFO records.}
S2 permits cancellation, while S4 selects a longest eligible FIFO prefix.
There is no deletion/tombstone rule. A canceled head could block later
entries under one reading, so the advertised FIFO progress is not yet
established. The completed request lifecycle needs explicit cleanup.

\item \textbf{Justification Gap: installation bookkeeping.}
R3 invokes S5 for departures but never explicitly promotes selected
\textsf{Reserved} joins to \textsf{Active}, consumes selected requests, or
clears the installed plan. EXIT requires an active position. These state
updates must be atomic parts of configuration installation.

\item \textbf{Justification Gap: waiting-branch precedence.}
R3 says to stay idle when the epoch is due but its plan is not ready;
R4 then appears to select a region unconditionally. The waiting branch
must explicitly bypass ordinary successor selection. Together with
V1-2, this leaves the advertised deterministic transition incomplete.
\end{enumerate}

\textbf{Verified dependency chain.} A finalized L1 expected-counter check
permits one transition per predecessor. Regional finality supplies prefix
compatibility inside a fixed session. The expressly assumed terminal
closure covers already-finalizable outputs and forbids later continuation.
Opening only the finalized successor at the closed anchor therefore
preserves prefixes by induction. The idle selector gap affects completeness
of the transition definition, not this induction once the intended branch
is fixed.

\textbf{Liability check.} Pending cancellation excludes reserved joins;
requested exits remain active; the cooling clock begins at actual
configuration retirement. Under the external timely-claim interface this
prevents voluntary withdrawal from erasing an unresolved timely claim.
It does not restore already slashed collateral.

\textbf{Accepted external premises.} L1 finality and ordering; sound
cryptography; fixed-configuration regional BFT safety with its separate
regional fault bounds; a valid terminal-command/closure adapter; data
availability for successor opening; the seed interface's uniqueness and
pre-freeze unpredictability; and correct escrow adjudication. The review
does not certify an implementation of these dependencies.

\section*{Patch plan for v2}
Specify the idle selector and no-op guard, make plan installation atomic,
remove canceled queue records, handle penalty-adjusted and zero balances,
limit churn claims to voluntary changes, and state configuration progress
separately from eligible-region production. Add explicit module API and
consumer receipt checks. Preserve the safe-waiting boundary on unavailable
closure and the separation of nominal deadlines from actual retirement.
\end{document}
